\loai{Tìm lực kéo vật lên}
\begin{vidu}%[1P7-3.2K][Loai1]
	Một vòng nhẫn rất nhẹ có đường kính trong là 4,5 cm và đường kính ngoài là 5 cm. Biết hệ số căng bề mặt ngoài của glyxêrin ở $20\doC$ là $65,2.10^{-3}\ N/m$. Tính lực bứt vòng nhẫn này ra khỏi mặt thoáng của glyxêrin?
	\choice
	{$9,2.10^{-3}\ N$}
	{$10,2.10^{-3}\ N$}
	{$6,1.10^{-3}\ N$}
	{\True $19,4.10^{-3}\ N$}
	\loigiai{
		\Binhluan{
		Lực bứt vòng nhẫn ra khỏi bề mặt đúng bằng lực căng bề mặt: 
		\begin{align*}
		f=\sigma.\ell=\sigma .\pi .(d+D)=19,4.10^{-3}\ N.
		\end{align*}
\begin{note}
	Bình thường lực căng sẽ nằm trong mặt phẳng nước; nhưng khi ta bứt vật ra khỏi mặt nước thì do hiện tượng dính ướt nên lực căng sẽ có phương hướng xuống.
\end{note}
}
{Vì vòng nhẫn rất nhẹ nên ta bỏ qua tác dụng của trọng lực P}
	}
\end{vidu}

\loai{Lực căng cân bằng với trọng lực}
\begin{vidu}%[1P7-3.2K][Loai3]
	Một ống nhỏ giọt dựng thẳng đứng bên trong đựng nước. Nước dính ướt hoàn toàn miệng ống và đường kính miệng dưới của ống là 0,43 mm. Trọng lượng mỗi giọt nước rơi khỏi miệng ống là $9,72.10^{-5}\ N$. Tính hệ số căng mặt ngoài của nước.
	\choice
	{$72.10^{-5}\ N/m$}
	{$36.10^{-3}\ N/m$}
	{\True $72.10^{-3}\ N/m$}
	{$13,8.10^2\ N/m$}
	\loigiai{
	\begin{itemize}
			\item Chiều dài đường giới hạn (chu vi đường tròn): $\ell=\pi d$.
			\item Lực căng mặt ngoài tác dụng lên đường giới hạn hướng thẳng đứng lên trên: 
			$f=\sigma.\ell=\sigma.d.\pi$
			\item Khi nước bắt đầu rơi khỏi miệng ống: 
			\begin{align*}
			f=P\Rightarrow \sigma.0,43.10^{-3}.3,14=9,72.10^{-5}\Rightarrow \sigma \approx 72.10^{-3}\ N/m.
			\end{align*}
		\end{itemize}
	}
\end{vidu}
\loai{Số giọt nước}
\begin{vidu}%[1P7-3.2K][Loai4]
	Cho $3\ cm^3$ nước vào ống nhỏ giọt đường kính 1 mm, thấy nhỏ được 120 giọt. Biết khối lượng riêng của nước là $D=1000\ kg/m^3$. Lấy $g=10\ m/s^2$. Tìm hệ số căng bề mặt của nước.
	\choice
	{\True 0,0795 N/m}
	{0,0242 N/m}
	{0,0213 N/m}
	{0,234 N/m}
	\loigiai{
		\Binhluan{
		\begin{itemize}
			\item Thể tích mỗi giọt là: $V_0=\dfrac{V}{n}=\dfrac{3.10^{-6}}{120}=2,5.10^{-8}\ m^2$.
			\item Khi nước bắt đầu nhỏ xuống: 
			\begin{align*}
			&f=P\Leftrightarrow \sigma.\pi d=V_0.D.g\\
			\Rightarrow\ &\sigma=\dfrac{V_0.D.g}{\pi d}=\dfrac{2,5.10^{-8}.1000.10}{\pi 0,001}= 0,0795\ N/m.
			\end{align*}
		\end{itemize}
	}
{Ống nhỏ giọt mỗi lần nhỏ ra một giọt có thể tích không đổi $V_0$}
	}
\end{vidu}
\begin{vidu}%[1P7-3.2G][Loai4]
	Một đơn thuốc có ghi: Mỗi ngày uống hai lần, mỗi lần 15 giọt. Biết suất căng mặt ngoài của thuốc là $8,5.10^{-2}\ N/m$, ống nhỏ giọt có đường kính 2 mm. Cho $g=10\ m/s^2$. Tính khối lượng thuốc uống mỗi ngày. 
	\choice
	{32 g}
	{8 g}
	{\True 16 g}
	{20 g}
	\loigiai{
		\begin{itemize}
			\item Ta coi rằng khi giọt thuốc rơi, trọng lượng giọt thuốc đúng bằng lực căng mặt ngoài tác dụng lên đường tròn giới hạn ở miệng ống $\Rightarrow P=f \Rightarrow mg=\pi d\sigma $.
			\item Vậy khối lượng một giọt nước là: 
			\begin{align*}
				m_1=\dfrac{{\pi d\sigma }}{g}=\dfrac{{3,14.8,{5.10}^{-2}.2.10}}{{10}}=53,38.10^{-3}\ kg.
			\end{align*}
			\item Khối lượng thuốc uống mỗi ngày là: $m=30.m_1=1,60.10^{-2} \ kg=16\ g$.
		\end{itemize}
	}
\end{vidu}


